%%%%% --------------------------------------------------------------------------------
%%
%%%%***************************** Literaure Review ***********************************
%%
%%% ++++++++++++++++++++++++++++++++++++++++++++++++++++++++++++++++++++++++++++++++++
\part{\FakeBold{文献综述}} % (fold)
\label{prt:文献综述_}
	\setcounter{section}{0}
	\section{背景介绍} % (fold)
	\label{sec:背景介绍}
		括号中斜体部分文字为说明文字。

		文献综述是研究者在阅读过某一主题的文献后，经过理解、整理、融会贯通，综合分析和评价而组成的一种不同于研究论文的文体。综述的目的是反映某一研究领域的新水平、新动态、新技术和新发现。一般需要对其从历史到现状、存在问题以及发展趋势等，进行全面的介绍和评论，为开题奠定良好的基础。

		文献综述的撰写遵照“文献综述与开题报告提纲”文件中给出的标题，内容可参考标题后括号内的文字，下级标题可根据综述的内容自行定义。

		文献综述的背景介绍可以包括与课题相关的研究背景，可以介绍有关的概念及定义以及综述的范围、扼要说明有关主题的现状，使读者对全文要叙述的问题有一个初步的轮廓。

		正文文字一级标题用小3号宋体加粗，正文采用小四号宋体，英文及数字采用Times New Roman，行间距为1.5倍行间距。字数要求3000字以上。

		如公式\eqref{eq:eqtest1}是单行公式环境，查看公式\eqref{eq:eqtest1}和\eqref{eq:eqtest2}之间的区别，好像在单行公式中没什么区别。
		\begin{equation}
		\label{eq:eqtest1}
			f(x) = 2(x + 1)^{2} - 1
		\end{equation}

		\texttt{align}公式环境，用在单行中。
		\begin{align}
		\label{eq:eqtest2}
			f(x) = 2(x + 1)^{2} - 1
		\end{align}
	% section 背景介绍 (end)
	\newpage
	\section{国内外研究现状} % (fold)
	\label{sec:国内外研究现状}
		国内外研究现状是文献综述的主体部分，可以包括历史发展、现状分析和研究展望等。在这部分，为了清晰地说明问题，可分为若干个子标题进行阐述。研究现状应将查阅的文献资料进行归纳、整理以及必要的分析，而不是简单的材料的堆砌。写法可以按年代顺序综述，也可按不同的问题进行综述，还可按不同的观点进行比较综述。无论采用方式，都需要做到全面系统、客观公正、分析透彻、层次分明、条理清楚、语言简练、详略得当。以下两个二级标题仅为说明格式要求。\citerev{bearman2015s}

		\subsection{研究方向及进展} % (fold)
		\label{sub:研究方向及进展}
			如\autoref{fig:image_tset}是对此模版的第一张插图测试。

			\begin{figure}[htbp]
				\centering
				\includegraphics[width = 0.5\linewidth]{./style/logo/ZJDX.pdf}
				\caption{\FakeBold{插图测试}}
				\label{fig:image_tset}
			\end{figure}
		% subsection 研究方向及进展 (end)
		\subsection{存在问题} % (fold)
		\label{sub:存在问题}
			以下节分别用来测试各种表格环境如，tabular，tabu，longtabu等，还有对caption格式的修改和测试。以下表格样式全部采用三线表。

			
			如\autoref{tab:first_table_test}是对array宏包的tabular表格环境测试。
			\begin{table}[htbp]
				\centering
				\zihao{5}
				\renewcommand{\arraystretch}{0.625}
				\caption{\FakeBold{这是一个用tabular环境的测试用的表格}}
				\label{tab:first_table_test}
			    \begin{tabular}{lrr}
			    \toprule
			    \textbf{行星}     & \textbf{赤道半径}km & \textbf{公转周期}d \\
			    \midrule
			    水星     & 2.439  & 87.9 \\
			    金星     & 6.1    & 224.682 \\
			    地球     & 6378.14 & 365.24 \\
			    \bottomrule
			    \end{tabular}%
			\end{table}

			
			如\autoref{tab:tabu_test_1}是对tabu宏包的tabu表格环境测试。在这里表格命令与\autoref{tab:first_table_test}的命令相同，只是tabular环境改成了tabu环境。
			\begin{table}[htbp]
				\centering
				\zihao{5}
				\renewcommand{\arraystretch}{0.625}
				\caption{\FakeBold{这是一个用tabu环境的测试用的表格}}
				\label{tab:tabu_test_1}
			    \begin{tabu}{lrr}
			    \toprule
			    \textbf{行星}     & \textbf{赤道半径}km & \textbf{公转周期}d \\
			    \midrule
			    水星     & 2.439  & 87.9 \\
			    金星     & 6.1    & 224.682 \\
			    地球     & 6378.14 & 365.24 \\
			    \bottomrule
			    \end{tabu}%
			\end{table}
		% subsection 存在问题 (end)
	% section 国内外研究现状 (end)
	\newpage
	\section*{参考文献} % (fold)
	\label{sec:参考文献1}
		\addcontentsline{toc}{section}{\protect\numberline{}参考文献}%
		\begingroup
		    \setlength{\bibsep}{0pt}
		    \setstretch{1}
		    \bibliographyrev{bib/literature_review}%
		\endgroup
	% section* 参考文献 (end)
% part 文献综述_ (end)

%%% ++++++++++++++++++++++++++++++++++++++++++++++++++++++++++++++++++++++++++++++++++
